\documentclass[12pt, a4paper]{ctexart}
\usepackage{fontspec}
\usepackage{xcolor}
\usepackage{hyperref}
\usepackage{geometry}
\usepackage{enumitem}
\usepackage{parskip}
\usepackage{titlesec}
\usepackage{tcolorbox}
\tcbuselibrary{breakable}
\usepackage{setspace}
\usepackage{listings}
\usepackage{amssymb}

\geometry{left=2.5cm, right=2.5cm, top=2.5cm, bottom=2.5cm}

\definecolor{myblue}{RGB}{30,100,200}
\definecolor{lightblue}{RGB}{230,240,255}
\definecolor{darkblue}{RGB}{0,50,120}

\newtcolorbox{bluebox}[1][]{
    colback=lightblue,
    colframe=myblue,
    coltitle=darkblue,
    fonttitle=\bfseries,
    title={#1},
    boxrule=0.8pt,
    arc=4pt,
    left=6pt,
    right=6pt,
    top=6pt,
    bottom=6pt,
    before skip=8pt,
    after skip=8pt,
    breakable
}

\usepackage{etoolbox}
\BeforeBeginEnvironment{bluebox}{\par\vfill}%
\AfterEndEnvironment{bluebox}{\par\vfill}%

\titleformat{\section}
  {\normalfont\Large\bfseries\color{myblue}}{\thesection}{1em}{}
\titlespacing*{\section}{0pt}{3.5ex plus 1ex minus .2ex}{1.5ex plus .2ex}

\lstset{
    language=Java,
    basicstyle=\ttfamily\small,
    keywordstyle=\color{blue},
    commentstyle=\color{green!60!black},
    stringstyle=\color{orange},
    numbers=left,
    numberstyle=\tiny\color{gray},
    frame=single,
    breaklines=true,
    columns=fullflexible,
    showstringspaces=false
}

\newcommand{\code}[1]{\texttt{#1}}

\newcounter{question}

\newenvironment{question}[1][]{%
    \refstepcounter{question}%
    \par\noindent\textbf{\thequestion. #1}%
    \par\vspace{0.4em}%
}{\par\vspace{0.8em}}

\newenvironment{options}{%
    \begin{enumerate}[label=\Alph*.]
}{\end{enumerate}}

\newenvironment{answer}{\par\noindent\textbf{答案:}\ }{\par}
\newenvironment{explanation}{\par\noindent\textbf{解析:}\ }{\par}

\begin{document}

\section{I/O}

\begin{question}
	以下哪个类用于实现文件读取？
	\begin{options}
		\item FileReader
		\item FileWriter
		\item BufferedReader
		\item BufferedWriter
	\end{options}
	\begin{answer}
		A
	\end{answer}
	\begin{explanation}
		\texttt{FileReader} 用于读取字符文件。
	\end{explanation}
\end{question}

\begin{bluebox}{答案与深度解析}
	\textbf{答案：A}

	% ======== 阶段一：核心认知框架 ========
	\textbf{【核心认知框架>}
	\begin{itemize}[label={}, leftmargin=*, nosep]
		\item \textbf{题目本质}：考查Java I/O流的基本分类与用途
		\item \textbf{关键区分点}：Reader/Writer处理字符流，InputStream/OutputStream处理字节流
		\item \textbf{命题人意图}：测试考生对Java I/O类库命名的理解
	\end{itemize}

	% ======== 阶段二：选项深度拆解 ========
	\begin{itemize}[label={}, nosep]
		\item \textbf{A. FileReader — 正确}
		\begin{itemize}[label={}, leftmargin=*, nosep]
			\item \textbf{继承关系}：
			      \begin{center}
				      \begin{tabular}{l}
					      \texttt{FileReader}        \\
					      $\downarrow$ extends       \\
					      \texttt{InputStreamReader} \\
					      $\downarrow$ extends       \\
					      \texttt{Reader}            \\
				      \end{tabular}
			      \end{center}

			\item \textbf{使用铁证}：
			      \begin{lstlisting}[language=Java]
try (FileReader reader = new FileReader("test.txt")) {
    int ch;
    while ((ch = reader.read()) != -1) {
        System.out.print((char) ch);
    }
}
      \end{lstlisting}

			\item \textbf{特点}：
			      \begin{itemize}[label={}, nosep]
				      \item 以\textbf{字符}为单位读取
				      \item 使用系统默认字符编码
				      \item 内部包装 FileInputStream
			      \end{itemize}

			\item \textbf{命题陷阱破解}：
			      \begin{itemize}[label={}, nosep]
				      \item 干扰项："\texttt{FileWriter} 也可以读" → 方向错误！
				      \item \textbf{必考结论}：\texttt{Reader} 系列用于读，\texttt{Writer} 系列用于写
			      \end{itemize}
		\end{itemize}

		\item \textbf{B. FileWriter — 错误（写入）}
		\begin{itemize}[label={}, leftmargin=*, nosep]
			\item \textbf{核心功能}：用于\textbf{写入}字符文件，不是读取

			\item \textbf{继承关系}：
			      \begin{center}
				      \begin{tabular}{l}
					      \texttt{FileWriter}         \\
					      $\downarrow$ extends        \\
					      \texttt{OutputStreamWriter} \\
					      $\downarrow$ extends        \\
					      \texttt{Writer}             \\
				      \end{tabular}
			      \end{center}

			\item \textbf{使用对比}：
			      \begin{lstlisting}[language=Java]
// FileWriter 用于写入
try (FileWriter writer = new FileWriter("output.txt")) {
    writer.write("Hello World");
}

// FileReader 用于读取
try (FileReader reader = new FileReader("output.txt")) {
    reader.read();
}
      \end{lstlisting}

			\item \textbf{命题陷阱破解}：
			      \begin{itemize}[label={}, nosep]
				      \item 干扰项："都带File肯定功能类似" → 只对了一半！
				      \item \textbf{必考结论}：看后缀 - \texttt{Reader}读，\texttt{Writer}写
			      \end{itemize}
		\end{itemize}

		\item \textbf{C. BufferedReader — 错误（需包装）}
		\begin{itemize}[label={}, leftmargin=*, nosep]
			\item \textbf{本质}：\textbf{装饰器类}，需要包装其他Reader
			      \begin{lstlisting}[language=Java]
// 错误用法：不能直接使用
BufferedReader reader = new BufferedReader("test.txt"); // 编译错误！

// 正确用法：包装FileReader
try (BufferedReader reader = new BufferedReader(
        new FileReader("test.txt"))) {
    String line = reader.readLine();
}
    \end{lstlisting}

			\item \textbf{BufferedReader特点}：
			      \begin{itemize}[label={}, nosep]
				      \item 提供\textbf{缓冲}功能，减少IO次数
				      \item 提供 \texttt{readLine()} 方法（常用！）
				      \item 默认8KB缓冲区
			      \end{itemize}

			\item \textbf{命题陷阱破解}：
			      \begin{itemize}[label={}, nosep]
				      \item 干扰项："Buffered系列更高级，应该用它" → 忽略需要包装！
				      \item \textbf{必考结论}：装饰器不能单独使用，必须包装底层流
			      \end{itemize}
		\end{itemize}

		\item \textbf{D. BufferedWriter — 错误（写入）}
		\begin{itemize}[label={}, leftmargin=*, nosep]
			\item \textbf{功能}：缓冲字符输出流，用于\textbf{写入}

			\item \textbf{使用铁证}：
			      \begin{lstlisting}[language=Java]
try (BufferedWriter writer = new BufferedWriter(
        new FileWriter("output.txt"))) {
    writer.write("Line 1");
    writer.newLine();
    writer.write("Line 2");
}
      \end{lstlisting}

			\item \textbf{与BufferedReader对比}：
			      \begin{center}
				      \begin{tabular}{|c|c|c|}
					      \hline
					      \textbf{类}     & \textbf{功能} & \textbf{特有方法} \\
					      \hline
					      BufferedReader & 读取          & readLine()    \\
					      \hline
					      BufferedWriter & 写入          & newLine()     \\
					      \hline
				      \end{tabular}
			      \end{center}

			\item \textbf{命题陷阱破解}：
			      \begin{itemize}[label={}, nosep]
				      \item 干扰项："Buffered比File更高级，应该用C/D" → 审题要看清"读"还是"写"
				      \item \textbf{必考结论}：先确定方向（读/写），再选具体类
			      \end{itemize}
		\end{itemize}
	\end{itemize}

	% ======== 阶段三：应背诵知识点 ========
	\textbf{【应背诵知识点（命题地人思维版）>}
	\begin{itemize}[label={}, nosep]
		\item \textbf{Java I/O分类}：
		      \begin{itemize}[label={}, nosep]
			      \item \textbf{字节流}：InputStream/OutputStream
			      \item \textbf{字符流}：Reader/Writer
			      \item \textbf{处理流}：Buffered、Print、Data等（需包装）
		      \end{itemize}

		\item \textbf{选项辨析速查表}：
		      \begin{center}
			      \begin{tabular}{|c|c|c|c|}
				      \hline
				      \textbf{选项}    & \textbf{功能} & \textbf{是否需包装} & \textbf{记忆口诀}    \\
				      \hline
				      FileReader     & 读取字符        & 否              & "Reader读"        \\
				      \hline
				      FileWriter     & 写入字符        & 否              & "Writer写"        \\
				      \hline
				      BufferedReader & 读取缓冲        & 需              & "Buffer包装Reader" \\
				      \hline
				      BufferedWriter & 写入缓冲        & 需              & "Buffer包装Writer" \\
				      \hline
			      \end{tabular}
		      \end{center}

		\item \textbf{高频陷阱清单}：
		      \begin{itemize}[label={}, nosep]
			      \item 陷阱1：装饰器类单独使用 → 编译错误
			      \item 陷阱2：只看前缀忽略后缀 → 读写混淆
			      \item 陷阱3：Buffered比File"更好" → 场景不同
		      \end{itemize}
	\end{itemize}

	% ======== 阶段四：命题趋势与实战策略 ========
	\textbf{【命题趋势与实战策略>}
	\begin{itemize}[label={}, leftmargin=*, nosep]
		\item \textbf{考场秒杀三步法}：
		      \begin{enumerate}
			      \item \textbf{先定方向}：题目问"读取" → Reader系列
			      \item \textbf{再看层级}：简单文件 → FileReader；需要缓冲 → BufferedReader包装
			      \item \textbf{避免陷阱}：Buffered系列必须包装
		      \end{enumerate}
	\end{itemize}
\end{bluebox}

\begin{question}
	以下哪个类用于实现文件写入？
	\begin{options}
		\item FileReader
		\item FileWriter
		\item BufferedReader
		\item BufferedWriter
	\end{options}
	\begin{answer}
		B
	\end{answer}
	\begin{explanation}
		\texttt{FileWriter} 用于写入字符文件。
	\end{explanation}
\end{question}

\begin{bluebox}{答案与深度解析}
	\textbf{答案：B}

	% ======== 阶段一：核心认知框架 ========
	\textbf{【核心认知框架>}
	\begin{itemize}[label={}, leftmargin=*, nosep]
		\item \textbf{题目本质}：考查Java文件写入类的识别
		\item \textbf{关键区分点}：FileWriter是写入文件的基础类
		\item \textbf{命题人意图}：测试考生对I/O类命名的理解
	\end{itemize}

	% ======== 阶段二：选项深度拆解 ========
	\begin{itemize}[label={}, nosep]
		\item \textbf{A. FileReader — 错误（读取）}
		\begin{itemize}[label={}, leftmargin=*, nosep]
			\item \textbf{用途}：用于读取，不是写入
			\item \textbf{命题陷阱}：带File但后缀是Reader，表示读取
		\end{itemize}

		\item \textbf{B. FileWriter — 正确}
		\begin{itemize}[label={}, leftmargin=*, nosep]
			\item \textbf{用途}：用于写入字符文件
			\item \textbf{使用铁证}：
			      \begin{lstlisting}[language=Java]
try (FileWriter writer = new FileWriter("test.txt")) {
    writer.write("Hello");
    writer.append(" World");
}
      \end{lstlisting}
		\end{itemize}

		\item \textbf{C. BufferedReader — 错误（读取+需包装）}
		\begin{itemize}[label={}, leftmargin=*, nosep]
			\item \textbf{用途}：读取缓冲流，需要包装FileReader
		\end{itemize}

		\item \textbf{D. BufferedWriter — 错误（需包装）}
		\begin{itemize}[label={}, leftmargin=*, nosep]
			\item \textbf{用途}：写入缓冲流，需要包装FileWriter
			\item \textbf{注意}：可以选，但题目问"哪个类"，FileWriter更基础
		\end{itemize}
	\end{itemize}

	% ======== 阶段三：应背诵知识点 ========
	\textbf{【应背诵知识点（命题人思维版）>}
	\begin{itemize}[label={}, nosep]
		\item \textbf{速记口诀}：Writer写，Reader读

		\item \textbf{选项辨析}：
		      \begin{center}
			      \begin{tabular}{|c|c|c|}
				      \hline
				      选项             & 功能   & 备注  \\
				      \hline
				      FileReader     & 读取   & 基础类 \\
				      \hline
				      FileWriter     & 写入   & 基础类 \\
				      \hline
				      BufferedReader & 读取缓冲 & 需包装 \\
				      \hline
				      BufferedWriter & 写入缓冲 & 需包装 \\
				      \hline
			      \end{tabular}
		      \end{center}
	\end{itemize}
\end{bluebox}

\begin{question}
	以下关于Java中的流的说法错误的是
	\begin{options}
		\item 分为字节流和字符流
		\item 字节流以字节为单位处理数据
		\item 字符流以字符为单位处理数据
		\item 字节流和字符流可以相互转换
	\end{options}
	\begin{answer}
		D
	\end{answer}
	\begin{explanation}
		字节流和字符流不能直接相互转换，需通过编码转换（如 \texttt{InputStreamReader}）。
	\end{explanation}
\end{question}

\begin{bluebox}{答案与深度解析}
	\textbf{答案：D}

	% ======== 阶段一：核心认知框架 ========
	\textbf{【核心认知框架>}
	\begin{itemize}[label={}, leftmargin=*, nosep]
		\item \textbf{题目本质}：考查Java I/O流的分类体系
		\item \textbf{关键区分点}：字节流和字符流是两种不同的处理方式，需要通过特定转换器转换
		\item \textbf{命题人意图}：测试考生对I/O流分类的掌握
	\end{itemize}

	% ======== 阶段二：选项深度拆解 ========
	\begin{itemize}[label={}, nosep]
		\item \textbf{A. 分为字节流和字符流 — 正确}
		\begin{itemize}[label={}, leftmargin=*, nosep]
			\item \textbf{分类体系}：
			      \begin{center}
				      \begin{tabular}{|c|c|c|}
					      \hline
					      \textbf{类型} & \textbf{基类}              & \textbf{处理单位}  \\
					      \hline
					      字节流         & InputStream/OutputStream & 字节（8位）         \\
					      \hline
					      字符流         & Reader/Writer            & 字符（16位Unicode） \\
					      \hline
				      \end{tabular}
			      \end{center}
		\end{itemize}

		\item \textbf{B. 字节流以字节为单位处理数据 — 正确}
		\begin{itemize}[label={}, leftmargin=*, nosep]
			\item \textbf{字节流铁证}：
			      \begin{lstlisting}[language=Java]
FileInputStream fis = new FileInputStream("test.txt");
int b = fis.read(); // 每次读取1个字节（0-255）
      \end{lstlisting}
		\end{itemize}

		\item \textbf{C. 字符流以字符为单位处理数据 — 正确}
		\begin{itemize}[label={}, leftmargin=*, nosep]
			\item \textbf{字符流铁证}：
			      \begin{lstlisting}[language=Java]
FileReader fr = new FileReader("test.txt");
int c = fr.read(); // 每次读取1个字符（Unicode）
      \end{lstlisting}
		\end{itemize}

		\item \textbf{D. 字节流和字符流可以相互转换 — 错误}
		\begin{itemize}[label={}, leftmargin=*, nosep]
			\item \textbf{转换铁证（间接转换）}：
			      \begin{lstlisting}[language=Java]
// 字节流转字符流：InputStreamReader
InputStream is = new FileInputStream("test.txt");
Reader reader = new InputStreamReader(is, "UTF-8");

// 字符流转字节流：OutputStreamWriter
Writer writer = new OutputStreamWriter(
    new FileOutputStream("test.txt"), "UTF-8");
      \end{lstlisting}

			\item \textbf{核心要点}：
			      \begin{itemize}[label={}, nosep]
				      \item \textbf{不能直接转换}：没有直接的字节流转字符流方法
				      \item \textbf{需要转换器}：InputStreamReader / OutputStreamWriter
				      \item \textbf{指定编码}：可以指定字符编码（UTF-8、GBK等）
			      \end{itemize}

			\item \textbf{命题陷阱破解}：
			      \begin{itemize}[label={}, nosep]
				      \item 干扰项："字符流就是字节流转换来的" → 忽略中间桥梁！
				      \item \textbf{必考结论}：需要\textbf{显式使用转换器}，不是自动转换
			      \end{itemize}
		\end{itemize}
	\end{itemize}

	% ======== 阶段三：应背诵知识点 ========
	\textbf{【应背诵知识点（命题人思维版）>}
	\begin{itemize}[label={}, nosep]
		\item \textbf{流分类速查表}：
		      \begin{center}
			      \begin{tabular}{|c|c|c|c|}
				      \hline
				      \textbf{流类型} & \textbf{基类}  & \textbf{单位} & \textbf{示例}      \\
				      \hline
				      字节输入流        & InputStream  & 字节          & FileInputStream  \\
				      \hline
				      字节输出流        & OutputStream & 字节          & FileOutputStream \\
				      \hline
				      字符输入流        & Reader       & 字符          & FileReader       \\
				      \hline
				      字符输出流        & Writer       & 字符          & FileWriter       \\
				      \hline
			      \end{tabular}
		      \end{center}

		\item \textbf{转换器}：
		      \begin{itemize}[label={}, nosep]
			      \item 字节→字符：\texttt{InputStreamReader}
			      \item 字符→字节：\texttt{OutputStreamWriter}
		      \end{itemize}

		\item \textbf{高频陷阱清单}：
		      \begin{itemize}[label={}, nosep]
			      \item 陷阱1：认为可以自动转换 → 需显式使用转换器
			      \item 陷阱2：忽略字符编码 → 转换时可指定编码
		      \end{itemize}
	\end{itemize}

	% ======== 阶段四：命题趋势与实战策略 ========
	\textbf{【命题趋势与实战策略】}
	\begin{itemize}[label={}, nosep]
		\item \textbf{考场秒杀三步法}：
		      \begin{enumerate}[label={}, nosep]
			      \item \textbf{回忆分类}：字节流 vs 字符流
			      \item \textbf{转换规则}：需要InputStreamReader/OutputStreamWriter
			      \item \textbf{选D}：D选项"可以直接转换"是错的
		      \end{enumerate}
	\end{itemize}
\end{bluebox}

\begin{question}
	以下哪个方法用于创建一个新文件？
	\begin{options}
		\item createNewFile()
		\item mkdir()
		\item mkdirs()
		\item exists()
	\end{options}
	\begin{answer}
		A
	\end{answer}
	\begin{explanation}
		\texttt{File.createNewFile()} 创建新文件。
	\end{explanation}
\end{question}

\begin{bluebox}{答案与深度解析}
	\textbf{答案：A}

	% ======== 阶段一：核心认知框架 ========
	\textbf{【核心认知框架>}
	\begin{itemize}[label={}, leftmargin=*, nosep]
		\item \textbf{题目本质}：考查File类的方法功能区分
		\item \textbf{关键区分点}：文件操作方法名相似但功能不同
		\item \textbf{命题人意图}：测试考生对File API的掌握
	\end{itemize}

	% ======== 阶段二：选项深度拆解 ========
	\begin{itemize}[label={}, nosep]
		\item \textbf{A. createNewFile() — 正确}
		\begin{itemize}[label={}, leftmargin=*, nosep]
			\item \textbf{功能}：创建新文件
			\item \textbf{铁证}：
			      \begin{lstlisting}[language=Java]
File file = new File("test.txt");
boolean created = file.createNewFile(); // 创建新文件
      \end{lstlisting}

			\item \textbf{特点}：
			      \begin{itemize}[label={}, nosep]
				      \item 若文件已存在，返回false
				      \item 若父目录不存在，会抛出IOException
			      \end{itemize}
		\end{itemize}

		\item \textbf{B. mkdir() — 错误（创建目录）}
		\begin{itemize}[label={}, leftmargin=*, nosep]
			\item \textbf{功能}：创建单层目录
			\item \textbf{区别}：创建目录而非文件
		\end{itemize}

		\item \textbf{C. mkdirs() — 错误（创建多层目录）}
		\begin{itemize}[label={}, leftmargin=*, nosep]
			\item \textbf{功能}：创建多层目录（含父目录）
		\end{itemize}

		\item \textbf{D. exists() — 错误（检查存在）}
		\begin{itemize}[label={}, leftmargin=*, nosep]
			\item \textbf{功能}：检查文件/目录是否存在
		\end{itemize}
	\end{itemize}

	% ======== 阶段三：应背诵知识点 ========
	\textbf{【应背诵知识点（命题人思维版）>}
	\begin{itemize}[label={}, nosep]
		\item \textbf{File方法速查}：
		      \begin{center}
			      \begin{tabular}{|c|c|c|}
				      \hline
				      \textbf{方法}   & \textbf{功能} & \textbf{返回值} \\
				      \hline
				      createNewFile & 创建文件        & boolean      \\
				      \hline
				      mkdir         & 创建单层目录      & boolean      \\
				      \hline
				      mkdirs        & 创建多层目录      & boolean      \\
				      \hline
				      exists        & 检查存在性       & boolean      \\
				      \hline
			      \end{tabular}
		      \end{center}

		\item \textbf{记忆口诀}：createNewFile创文件，mkdir创单层，mkdirs创多层
	\end{itemize}
\end{bluebox}

\begin{question}
	以下哪个方法用于删除一个文件？
	\begin{options}
		\item delete()
		\item remove()
		\item clear()
		\item erase()
	\end{options}
	\begin{answer}
		A
	\end{answer}
	\begin{explanation}
		\texttt{File.delete()} 删除文件或目录。
	\end{explanation}
\end{question}

\begin{bluebox}{答案与深度解析}
	\textbf{答案：A}

	% ======== 阶段一：核心认知框架 ========
	\textbf{【核心认知框架>}
	\begin{itemize}[label={}, leftmargin=*, nosep]
		\item \textbf{题目本质}：考查File类的删除方法
		\item \textbf{关键区分点}：File类只有delete方法，没有remove/clear/erase
		\item \textbf{命题人意图}：测试考生对File API的掌握
	\end{itemize}

	% ======== 阶段二：选项深度拆解 ========
	\begin{itemize}[label={}, nosep]
		\item \textbf{A. delete() — 正确}
		\begin{itemize}[label={}, leftmargin=*, nosep]
			\item \textbf{功能}：删除文件或空目录
			\item \textbf{铁证}：
			      \begin{lstlisting}[language=Java]
File file = new File("test.txt");
boolean deleted = file.delete(); // 删除文件
      \end{lstlisting}
		\end{itemize}

		\item \textbf{B/C/D — 错误}
		\begin{itemize}[label={}, leftmargin=*, nosep]
			\item \textbf{注意}：这些方法在File类中不存在
		\end{itemize}
	\end{itemize}

	% ======== 阶段三：应背诵知识点 ========
	\textbf{【应背诵知识点（命题人思维版）>}
	\begin{itemize}[label={}, nosep]
		\item \textbf{速记}：File类删除只有delete()方法

		\item \textbf{delete()特点}：
		      \begin{itemize}[label={}, nosep]
			      \item 文件：直接删除
			      \item 目录：仅当目录为空时才能删除
		      \end{itemize}
	\end{itemize}
\end{bluebox}

\begin{question}
	以下关于Java中序列化的说法错误的是
	\begin{options}
		\item 可以将对象转换为字节序列
		\item 可以通过实现Serializable接口实现序列化
		\item 静态成员变量不能被序列化
		\item 所有成员变量都会被序列化
	\end{options}
	\begin{answer}
		D
	\end{answer}
	\begin{explanation}
		被 \texttt{transient} 修饰的变量不会被序列化。
	\end{explanation}
\end{question}

\begin{bluebox}{答案与深度解析}
	\textbf{答案：D}

	% ======== 阶段一：核心认知框架 ========
	\textbf{【核心认知框架>}
	\begin{itemize}[label={}, leftmargin=*, nosep]
		\item \textbf{题目本质}：考查Java序列化机制的细节
		\item \textbf{关键区分点}：序列化有例外情况（static、transient）
		\item \textbf{命题人意图}：测试考生对序列化边界条件的理解
	\end{itemize}

	% ======== 阶段二：选项深度拆解 ========
	\begin{itemize}[label={}, nosep]
		\item \textbf{A. 可以将对象转换为字节序列 — 正确}
		\begin{itemize}[label={}, leftmargin=*, nosep]
			\item \textbf{序列化定义}：对象 $\rightarrow$ 字节序列
			\item \textbf{反序列化}：字节序列 $\rightarrow$ 对象
		\end{itemize}

		\item \textbf{B. 可以通过实现Serializable接口实现序列化 — 正确}
		\begin{itemize}[label={}, leftmargin=*, nosep]
			\item \textbf{标记接口}：Serializable没有方法，只起标记作用
			\item \textbf{铁证}：
			      \begin{lstlisting}[language=Java]
public class Person implements Serializable {
    private static final long serialVersionUID = 1L;
    private String name;
    private int age;
}
      \end{lstlisting}
		\end{itemize}

		\item \textbf{C. 静态成员变量不能被序列化 — 正确}
		\begin{itemize}[label={}, leftmargin=*, nosep]
			\item \textbf{原因}：静态变量属于类，不属于对象
			\item \textbf{铁证}：
			      \begin{lstlisting}[language=Java]
public class Person implements Serializable {
    static int count = 0; // 静态变量

    public static void main(String[] args) throws Exception {
        Person p = new Person();
        p.count = 100;

        // 序列化
        ByteArrayOutputStream baos = new ByteArrayOutputStream();
        ObjectOutputStream oos = new ObjectOutputStream(baos);
        oos.writeObject(p);
        oos.close();

        // 反序列化
        ByteArrayInputStream bais = new ByteArrayInputStream(baos.toByteArray());
        ObjectInputStream ois = new ObjectInputStream(bais);
        Person p2 = (Person) ois.readObject();

        System.out.println(p2.count); // 输出0，非100！static未被序列化
    }
}
      \end{lstlisting}
		\end{itemize}

		\item \textbf{D. 所有成员变量都会被序列化 — 错误}
		\begin{itemize}[label={}, leftmargin=*, nosep]
			\item \textbf{例外1：transient}：
			      \begin{lstlisting}[language=Java]
public class Person implements Serializable {
    String name;
    transient String password; // 不会被序列化

    public static void main(String[] args) throws Exception {
        Person p = new Person();
        p.name = "Tom";
        p.password = "secret";

        // 序列化后反序列化
        // password会变成null
    }
}
      \end{lstlisting}

			\item \textbf{例外2：static}：见选项C说明

			\item \textbf{例外3：非Serializable对象}：
			      \begin{itemize}[label={}, nosep]
				      \item 成员变量类型未实现Serializable
				      \item 序列化时会抛出NotSerializableException
			      \end{itemize}

			\item \textbf{命题陷阱破解}：
			      \begin{itemize}[label={}, nosep]
				      \item 干扰项："所有字段都会被保存" → 忽略例外情况！
				      \item \textbf{必考结论}：\textbf{static + transient} 不参与序列化
			      \end{itemize}
		\end{itemize}
	\end{itemize}

	% ======== 阶段三：应背诵知识点 ========
	\textbf{【应背诵知识点（命题人思维版）>}
	\begin{itemize}[label={}, nosep]
		\item \textbf{序列化核心规则}：
		      \begin{center}
			      \begin{tabular}{|c|c|c|}
				      \hline
				      \textbf{类型}         & \textbf{是否序列化} & \textbf{原因} \\
				      \hline
				      普通字段                & \checkmark     & 默认行为        \\
				      \hline
				      static字段            & \textbf{X}     & 属于类非对象      \\
				      \hline
				      transient字段         & \textbf{X}     & 显式排除        \\
				      \hline
				      实现Serializable的对象字段 & \checkmark     & 递归序列化       \\
				      \hline
			      \end{tabular}
		      \end{center}

		\item \textbf{选项辨析速查表}：
		      \begin{tabular}{|c|c|c|c|}
			      \hline
			      \textbf{选项}    & \textbf{正误} & \textbf{记忆口诀} \\
			      \hline
			      转字节序列          & \checkmark  & 序列化本质         \\
			      \hline
			      实现Serializable & \checkmark  & 标记接口          \\
			      \hline
			      static不序列化     & \checkmark  & 属于类           \\
			      \hline
			      所有都序列化         & \textbf{X}  & transient排除   \\
			      \hline
		      \end{tabular}

		\item \textbf{高频陷阱清单}：
		      \begin{itemize}[label={}, nosep]
			      \item 陷阱1：static被序列化 → 错！属于类
			      \item 陷阱2：transient被序列化 → 错！显式排除
			      \item 陷阱3：所有字段都保存 → 错！有例外
		      \end{itemize}
	\end{itemize}

	% ======== 阶段四：命题趋势与实战策略 ========
	\textbf{【命题趋势与实战策略>}
	\begin{itemize}[label={}, leftmargin=*, nosep]
		\item \textbf{考场秒杀三步法}：
		      \begin{enumerate}[label={}, nosep]
			      \item \textbf{记住两个例外}：static + transient
			      \item \textbf{Serializable接口}：标记接口
			      \item \textbf{serialVersionUID}：建议显式声明
		      \end{enumerate}

		\item \textbf{高阶延伸}：
		      \begin{itemize}[label={}, nosep]
			      \item Externalizable：自定义序列化
			      \item writeObject/readObject：自定义序列化逻辑
			      \item JSON/XML替代方案：跨语言/可读性好
		      \end{itemize}
	\end{itemize}
\end{bluebox}

\begin{question}
	以下哪个方法用于反序列化？
	\begin{options}
		\item readObject()
		\item deserialize()
		\item restoreObject()
		\item recoverObject()
	\end{options}
	\begin{answer}
		A
	\end{answer}
	\begin{explanation}
		\texttt{ObjectInputStream.readObject()} 用于反序列化。
	\end{explanation}
\end{question}

\begin{bluebox}{答案与深度解析}
	\textbf{答案：A}

	% ======== 阶段一：核心认知框架 ========
	\textbf{【核心认知框架>}
	\begin{itemize}[label={}, leftmargin=*, nosep]
		\item \textbf{题目本质}：考查Java反序列化API
		\item \textbf{关键区分点}：只有readObject()是标准API
		\item \textbf{命题人意图}：测试考生对序列化API的掌握
	\end{itemize}

	% ======== 阶段二：选项深度拆解 ========
	\begin{itemize}[label={}, nosep]
		\item \textbf{A. readObject() — 正确}
		\begin{itemize}[label={}, leftmargin=*, nosep]
			\item \textbf{来源}：ObjectInputStream类
			\item \textbf{铁证}：
			      \begin{lstlisting}[language=Java]
FileInputStream fis = new FileInputStream("person.ser");
ObjectInputStream ois = new ObjectInputStream(fis);
Person p = (Person) ois.readObject(); // 反序列化
ois.close();
      \end{lstlisting}
		\end{itemize}

		\item \textbf{B/C/D — 错误}
		\begin{itemize}[label={}, leftmargin=*, nosep]
			\item \textbf{注意}：这些方法不是Java标准API中的方法
		\end{itemize}
	\end{itemize}

	% ======== 阶段三：应背诵知识点 ========
	\textbf{【应背诵知识点（命题人思维版）>}
	\begin{itemize}[label={}, nosep]
		\item \textbf{序列化/反序列化API}：
		      \begin{center}
			      \begin{tabular}{|c|c|c|}
				      \hline
				      \textbf{操作} & \textbf{类}         & \textbf{方法}   \\
				      \hline
				      序列化         & ObjectOutputStream & writeObject() \\
				      \hline
				      反序列化        & ObjectInputStream  & readObject()  \\
				      \hline
			      \end{tabular}
		      \end{center}

		\item \textbf{速记口诀}：序列化用write，反序列化用read
	\end{itemize}
\end{bluebox}

\end{document}
